% ============================================================
%  Requisitos del sistema
% ============================================================
\section{Requisitos del sistema a desarrollar}
\label{sec:requisitos}

\subsection{Hardware}
\label{subsec:requisitos_hw}

El subsistema de hardware constituye el núcleo físico del equipo. Su función es
ejecutar, de forma controlada y repetible, los procedimientos de medición
definidos por la norma IEC~60335, garantizando la seguridad del operario y la
trazabilidad de los resultados. El dispositivo debe cubrir:
\begin{itemize}
    \item Medición de consumo (tensión, corriente, potencia activa y factor de
          potencia).
    \item Medición de corriente de fuga activa y pasiva.
    \item Medición de resistencia de puesta a tierra.
    \item Ensayo de rigidez dieléctrica (tensión resistida).
    \item Conmutación automática entre los circuitos de ensayo y desconexión
          del aparato bajo ensayo (DUT).
    \item Medidas de seguridad: aislamiento entre etapas, protecciones,
          desconexión ante falla, indicadores y paro de emergencia.
\end{itemize}

\subsection{Software}
\label{subsec:requisitos_sw}

La interfaz debe abarcar:
\begin{itemize}
    \item Ejecución de los cinco ensayos, adaptando los límites a los atributos
          del producto (clase de aislamiento, tipo de aparato, cable de
          alimentación, condición de duplicación del límite de fuga).
    \item Configuración de productos y tiempos de ensayo, secuencia automática de
          ensayos, adquisición y procesamiento de datos y veredicto automático.
    \item Alarmas ante valores anormales o fallas del equipo.
    \item Registro histórico trazable y exportación de resultados.
\end{itemize}
